\documentclass[10pt,a4paper]{amsart}
\usepackage[margin=19mm]{geometry}
\usepackage{amsmath,amssymb,mathtools}
\usepackage[scr=rsfs,cal=euler]{mathalfa}
\usepackage{xfrac}
\usepackage[unicode,hidelinks,bookmarks=false]{hyperref}
\newcommand{\ie}{i.e.\ }
\newcommand{\cf}{cf.\ }
\newcommand{\resp}{respectively\ }
\newcommand{\Subsection}{\subsection}
\providecommand{\nobreakdash}{\nobreak\hbox{-}}
\setlength{\emergencystretch}{2em}
\multlinegap0pt
\usepackage{amssymb}
\newtheorem{proposition}{Proposition}
\renewcommand{\le}{\leqslant}
\title{On $\eta$-periodic Formal Ternary Laws}
\author{Tao Huang}
\date{Source: arXiv:2607.06795v2 (CC BY 4.0)}
\begin{document}
\maketitle

\noindent\emph{Source and licence.} On $\eta$-periodic Formal Ternary Laws. Version arXiv:2607.06795v2 (CC BY 4.0), distributed by arXiv under the Creative Commons Attribution 4.0 International licence, http://creativecommons.org/licenses/by/4.0/, as recorded in the arXiv licence metadata for this exact version and shown on the abstract page https://arxiv.org/abs/2607.06795v2. Full licence notice: \texttt{licenses/arxiv-2607.06795-CC-BY-4.0.txt}.\par\smallskip

\noindent\textit{Editorial scope.} Counted unit: the article's proposition prop:ftl-nonpositive-coefficients (source0.tex lines 1394--1396). The article's complete original proof of it is delivered with it (source0.tex lines 1397--1459, one \texttt{proof} environment of 63 lines). No mathematics is written, repaired or re-derived: the statement and the proof are the article's own lines, in the article's own notation, and no retained line is substituted or re-flowed.  No further source-local context is needed: every cross-reference of every retained line resolves inside the two delivered spans.  Nothing was omitted from any retained span, so the package carries no omission record.  No retained line contains a citation, so the wrapper carries no bibliography and no bibliography entry.  No retained line contains a figure environment, an image-inclusion command or drawing code, so no image is delivered and the package declares no asset dependency.  Editorial formatting only, in addition: an AMS \texttt{amsart} preamble that restates the article's own declarations and macros used by the retained lines, namely the environments \texttt{proposition} on separate counters, together with the standard packages those lines need.  The retained environments print in this document as Proposition 1 in order of appearance, whereas the article prints the same items with the numbering of its own full document.  The complete native source of the article as distributed in the arXiv e-print for this exact version is bundled unchanged as \texttt{sources/204661/source0.tex}.
\begin{proposition}\label{prop:ftl-nonpositive-coefficients}
	The non-positive degree coefficients of $F_t(x,y,z)$ are completely determined by the axioms.
\end{proposition}

\begin{proof}
	First replace $F_t$ by the normalized law $\Psi_*F_t$, and write it again as $F_t$. Its $\iota$-linearity gives
	\[
		F_l^{\le 0}(-x,y,z)=(-1)^l F_l^{\le 0}(x,y,z).
	\]
	Combined with symmetry and $2$-saturation, each monomial $x^ay^bz^c$ in $F_l^{\le 0}$ must have all exponents congruent to $l$ modulo~$2$. In particular $F_1^{\le 0}=0$, and
	\[
		\begin{aligned}
			            & F_2^{\le 0}=a^2_{2,0,0}\sigma(x^2),\qquad
			F_3^{\le 0}=a^3_{1,1,1}xyz,                                                             \\
			F_4^{\le 0} & =a^4_{2,0,0}\sigma(x^2)+a^4_{4,0,0}\sigma(x^4)+a^4_{2,2,0}\sigma(x^2y^2).
		\end{aligned}
	\]

	From the neutral element axiom
	\[
		F_t(x,0,0)=(1+xt)^2(1+\iota(x)t)^2
		=1-2x^2t^2+x^4t^4+\text{(terms of higher $x$-degree)},
	\]
	we obtain
	\[
		a^2_{2,0,0}=-2,\qquad
		a^4_{2,0,0}=0,\qquad
		a^4_{4,0,0}=1.
	\]
	Next, from the weak neutral element axiom
	\[
		F_t(x,x,y)=(1+yt)^2(1+S(x,y)t+P(x,y)t^2),
	\]
	the $t$- and $t^2$-coefficients give
	\[
		F_1(x,x,y)=2y+S(x,y),\qquad
		F_2(x,x,y)=y^2+2yS(x,y)+P(x,y).
	\]
	Since $F_1^{\le 0}=0$ and $F_2^{\le 0}(x,y,z)=-2\sigma(x^2)$, we obtain
	\[
		S^{\le 0}(x,y)=-2y,\qquad
		P^{\le 0}(x,y)=-4x^2+y^2.
	\]
	The $t^3$-coefficient is
	\[
		F_3(x,x,y)=2yP(x,y)+y^2S(x,y).
	\]
	Comparing the coefficient of $x^2y$, only the term $2yP(x,y)$ contributes, so
	\[
		a^3_{1,1,1}=[F_3(x,x,y)]_{x^2y}=2[P(x,y)]_{x^2}=2(-4)=-8.
	\]
	Hence
	\[
		F_2^{\le 0}(x,y,z)=-2\sigma(x^2),\qquad
		F_3^{\le 0}(x,y,z)=-8xyz.
	\]

	It remains to determine $a^4_{2,2,0}$. From the weak neutral element axiom, the $t^4$-coefficient satisfies $F_4(x,x,y)=y^2P(x,y)$. Using the formula for $P^{\le 0}$ above and comparing the degree-$4$ part of $F_4^{\le 0}(x,x,y)=y^2P^{\le 0}(x,y)$:
	\[
		(2+a^4_{2,2,0})x^4+2a^4_{2,2,0}x^2y^2+y^4=-4x^2y^2+y^4,
	\]
	giving $a^4_{2,2,0}=-2$. Therefore
	\[
		F_4^{\le 0}(x,y,z)=\sigma(x^4)-2\sigma(x^2y^2).
	\]
	This also shows that the normalized law has no negative-degree coefficients. Since the nonlinear coefficients of $\Psi^{\pm1}$ have positive degree, transporting back preserves the displayed non-positive part.
\end{proof}

\end{document}
